% !TEX root = ../main.tex
\section{Hard Instances for the Lower Bounds}
\label{app:lower-bound}

This appendix gives the proofs deferred from Section~\ref{sec:lb-proof}
and the complete analyses of the hard instances. Throughout, we set
\(\mathbf C_0:=(1,1,0,1,1)\). Since \(\mathcal I(\mathbf C,\Delta)\) and
\(\mathcal S_0(\mathbf C)\) are monotone in \(\mathbf C\), membership for
\(\mathbf C_0\) implies membership for every \(\mathbf C\) with
\(L_f,G_f,G_{\calA},L_h\ge1\).

\subsection{Reduction to a gradient descent-ascent recursion}
\label{subsec:gda-embedding}
\label{app:gda-embedding-proof}

Define the region
\(\mathcal R_\rho:=\{(\bx,\blmbd)\in\RR^n\times\RR^m:\blmbd+\rho\calA(\bx)\in\partial h(\zero)\}\)
and the smooth nonconvex-linear function
\[
    \Psi_\rho(\bx,\blmbd):=f(\bx)+\frac{\rho}{2}\|\calA(\bx)\|^2
    +\langle\blmbd,\calA(\bx)\rangle.
\]

\begin{proposition}[Local linearity of \(\calD_\rho\) on \(\mathcal R_\rho\)]
\label{prop:reduction}
Let \(\rho>0\) and \((\bx,\blmbd)\in\mathcal R_\rho\). Then
\(\by(\bx,\blmbd)=\zero\), \(\bq(\bx,\blmbd)=\blmbd+\rho\calA(\bx)\),
\[
    \nabla_{\bx}\calD_\rho(\bx,\blmbd)=\nabla_{\bx}\Psi_\rho(\bx,\blmbd),
    \qquad
    \nabla_{\blmbd}\calD_\rho(\bx,\blmbd)=\nabla_{\blmbd}\Psi_\rho(\bx,\blmbd)
    =\calA(\bx),
\]
and \(\calD_\rho(\bx,\blmbd)=\Psi_\rho(\bx,\blmbd)+h(\zero)\).
\end{proposition}

\begin{proof}
Set \(\bv:=\calA(\bx)+\blmbd/\rho\). By~\eqref{eq:prox-multiplier}, \(\by(\bx,\blmbd)\)
is the unique minimizer of \(h(\by)+\frac{\rho}{2}\|\by-\bv\|^2\), and
\(\by\) is this minimizer if and only if \(\rho(\bv-\by)\in\partial h(\by)\).
For \(\by=\zero\) the inclusion reads
\(\blmbd+\rho\calA(\bx)\in\partial h(\zero)\), which holds because
\((\bx,\blmbd)\in\mathcal R_\rho\). Hence \(\by(\bx,\blmbd)=\zero\) and
\(\bq(\bx,\blmbd)=\blmbd+\rho\calA(\bx)\). Since
\(\nabla_{\bx}\bigl(\frac{\rho}{2}\|\calA(\bx)\|^2\bigr)=\rho J_{\calA}(\bx)^\top\calA(\bx)\),
\eqref{eq:reduced-al-gradient} gives
\[
\nabla_{\bx}\calD_\rho(\bx,\blmbd)=
\nabla f(\bx)+J_{\calA}(\bx)^\top\bigl(\blmbd+\rho\calA(\bx)\bigr)
=\nabla_{\bx}\Psi_\rho(\bx,\blmbd),
\]
and \(\nabla_{\blmbd}\calD_\rho(\bx,\blmbd)=\calA(\bx)-\zero=\calA(\bx)=\nabla_{\blmbd}\Psi_\rho(\bx,\blmbd)\).
Finally, completing the square in~\eqref{eq:reduced-al} gives
\(\calD_\rho(\bx,\blmbd)=f(\bx)+\min_{\by}\{h(\by)+\frac{\rho}{2}\|\by-\bv\|^2\}
-\frac{1}{2\rho}\|\blmbd\|^2\), and the minimum equals
\(h(\zero)+\frac{\rho}{2}\|\bv\|^2\) because its minimizer is \(\zero\). Hence
\[
\calD_\rho(\bx,\blmbd)
=f(\bx)+h(\zero)+\frac{\rho}{2}\|\calA(\bx)\|^2+\langle\blmbd,\calA(\bx)\rangle
+\frac{1}{2\rho}\|\blmbd\|^2-\frac{1}{2\rho}\|\blmbd\|^2
=\Psi_\rho(\bx,\blmbd)+h(\zero).
\]
\end{proof}

For \(\theta=(\rho,s,\gamma,\eta,\tau)\) and \((\alpha,\beta)\in[0,1]^2\), consider
the recursion
\begin{equation}\label{eq:gda-recursion}
\begin{aligned}
\bx^{k+1}&=\bx^k-\eta\Bigl[\nabla_{\bx}\Psi_\rho(\bx^k,\blmbd^k)
  +\frac{1}{\gamma}(\bx^k-\bz^k)\Bigr],\\
\blmbd^{k+1}&=\blmbd^k+\tau\Bigl[\nabla_{\blmbd}\Psi_\rho(\bx^{k+1},\blmbd^k)
  -\frac{1}{s}(\blmbd^k-\bmu^k)\Bigr],\\
\bz^{k+1}&=\bz^k+\alpha(\bx^{k+1}-\bz^k),\qquad
\bmu^{k+1}=\bmu^k+\beta(\blmbd^{k+1}-\bmu^k).
\end{aligned}
\end{equation}

\begin{proposition}[Embedding]
\label{prop:embedding}
Let \(\rho>0\), and let the sequence
\(\{(\bx^k,\blmbd^k,\bz^k,\bmu^k)\}_{k\ge0}\) be generated by
\(\mathrm{LRALM}_{\theta,\alpha,\beta}\). Suppose that for some \(K\ge1\),
\((\bx^k,\blmbd^k)\in\mathcal R_\rho\) and
\((\bx^{k+1},\blmbd^k)\in\mathcal R_\rho\) for all \(0\le k<K\).
Then~\eqref{eq:gda-recursion} holds for \(0\le k<K\), and
\[
R_k=\max\bigl\{\|\nabla_{\bx}\Psi_\rho(\bx^{k+1},\blmbd^k)\|,\;
\|\nabla_{\blmbd}\Psi_\rho(\bx^{k+1},\blmbd^k)\|\bigr\},
\qquad 0\le k<K.
\]
\end{proposition}

\begin{proof}
Apply Proposition~\ref{prop:reduction} at \((\bx^k,\blmbd^k)\) in the
\(\bx\)-update of~\eqref{eq:splalm-updates}, and at \((\bx^{k+1},\blmbd^k)\) in
the \(\blmbd\)-update. The same point gives \(\by^{k+1}=\zero\) and
\(\bq^k=\blmbd^k+\rho\calA(\bx^{k+1})\), so the formula for \(R_k\) follows
from~\eqref{eq:kkt-residual}.
\end{proof}


\subsection{Proof of Lemma~\ref{lem:main-plateau}}
\label{subsec:lb-plateau}

\begin{proof}[Proof of Lemma~\ref{lem:main-plateau}]
A direct check shows that \((f,\calA,h)\in\mathcal I(\mathbf C_0,\Delta)\) with
\(f(0)=\Delta\) and \(\inf f=f(x_\star)=0\).
We establish three properties of the iterates of
\(\mathrm{LRALM}_{\theta,\alpha,\beta}\) on Example~\ref{ex:plateau} from
\(\bu^0=\zero\).

(i) Since \(\calA\equiv0\) and \(\partial h(0)=[-1,1]\), every pair
\((x,0)\) lies in \(\mathcal R_\rho\), and \(\Psi_\rho(x,\lambda)=f(x)\). The
\(\lambda\)-update in~\eqref{eq:gda-recursion} reads
\(\lambda^{k+1}=\lambda^k-\frac{\tau}{s}(\lambda^k-\mu^k)\). Starting from
\(\lambda^0=\mu^0=0\), induction gives \(\lambda^k=\mu^k=0\) for all \(k\).
Proposition~\ref{prop:embedding} then gives \(y^{k+1}=0\) and
\(R_k=\max\{\|\nabla f(x^{k+1})\|,0\}=\|\nabla f(x^{k+1})\|\).

(ii) We use the Euclidean norm, which coincides with the absolute value
in one dimension. By (i), the primal update gives
\[
x^{k+1}-z^k
=\left(1-\frac{\eta}{\gamma}\right)(x^k-z^k)-\eta\nabla f(x^k).
\]
Since \(\|\nabla f(x)\|\le2\epsilon\) for every \(x\) and
\(0<\eta\le\gamma\),
\[
\|x^{k+1}-z^k\|
\le\left(1-\frac{\eta}{\gamma}\right)\|x^k-z^k\|+2\eta\epsilon.
\]
The center update implies
\[
\|x^{k+1}-z^{k+1}\|=(1-\alpha)\|x^{k+1}-z^k\|,
\qquad
\|z^{k+1}-z^k\|=\alpha\|x^{k+1}-z^k\|.
\]
Starting from \(x^0=z^0=0\), induction gives
\(\|x^k-z^k\|\le2\gamma\epsilon\) and
\(\|x^{k+1}-z^k\|\le2\gamma\epsilon\) for all \(k\ge0\).
Consequently,
\[
\|z^k-z^0\|
\le\sum_{j=0}^{k-1}\|z^{j+1}-z^j\|
\le2\alpha\gamma\epsilon k,
\qquad
\|x^k-x^0\|
\le\|x^k-z^k\|+\|z^k-z^0\|
\le2\gamma\epsilon(1+\alpha k).
\]
If \(\alpha>0\), then for every integer \(k\le\Delta/(16\gamma\alpha\epsilon^2)\), the
assumptions \(\epsilon^2\le\Delta/(16\gamma)\) and
\(\epsilon^2\le\Delta/4\) imply
\[
\|x^k-x^0\|
\le2\gamma\epsilon+\frac{\Delta}{8\epsilon}
\le\frac{\Delta}{4\epsilon}
\le\frac{\Delta}{2\epsilon}-\epsilon
=x_\star-2\epsilon.
\]
In particular, \(x^k\le\|x^k-x^0\|\le x_\star-2\epsilon\), so the
iterate lies in the constant-gradient region. When \(\alpha=0\), the
center remains at \(z^k=z^0=0\), and the same conclusion holds for every
\(k\ge0\).

(iii) Suppose \(\alpha>0\). By (ii), every output with
\(k+1\le\Delta/(16\gamma\alpha\epsilon^2)\) satisfies
\(x^{k+1}\le x_\star-2\epsilon\). By (i), its residual is
\(R_k=\|\nabla f(x^{k+1})\|=2\epsilon>\epsilon\).
Thus, none of the first
\(\lfloor\Delta/(16\gamma\alpha\epsilon^2)\rfloor\) outputs is an
\(\epsilon\)-KKT point. Definition~\ref{def:complexity-measure} yields
\[
T_\epsilon(\theta,\alpha,\beta;I,\zero)
\ge\left\lfloor\frac{\Delta}{16\gamma\alpha\epsilon^2}\right\rfloor+1
>\frac{\Delta}{16\gamma\alpha\epsilon^2}.
\]
When \(\alpha=0\), (ii) keeps every iterate in the constant-gradient region.
Hence \(R_k=2\epsilon>\epsilon\) for all \(k\ge0\), and
\(T_\epsilon(\theta,\alpha,\beta;I,\zero)=+\infty\).
\end{proof}

\subsection{Proof of Lemma~\ref{lem:main-rotation}}
\label{subsec:lb-rotation}

We analyze \(\mathrm{LRALM}_{\theta,\alpha,\beta}\) on Example~\ref{ex:rotation}
initialized at \(\bu^0=\zero\). Since \(h=\|\cdot\|_2\), Proposition~\ref{prop:reduction} ensures that both proximal outputs vanish
whenever
\[
\max_{j\in\{k,k+1\}}\|\lambda^k+\rho\calA(x^j)\|\le1.
\] 
Under this condition, LRALM follows~\eqref{eq:gda-recursion}. In the regime where \(\|x-x_c\|\le R_*\) and \(\|\lambda+\rho\calA(x)\|\le1\),
the reduced function takes the quadratic form
\[
\Psi_\rho(x,\lambda)
=-\frac{a}{2}\|x-x_c\|^2+b\langle\lambda-\lambda_c,x-x_c\rangle .
\]
Here \(a:=c_\kappa-\rho b^2\). The gradients of this quadratic function are
linear in \(x-x_c\) and \(\lambda-\lambda_c\). 
Therefore, let define the iteration sequence \(\bu^k:=(x^k,\lambda^k,z^k,\mu^k)\) and
\(\bu_c:=(x_c,\lambda_c,x_c,\lambda_c)\). Whenever the proximal condition
above and \(\|x^k-x_c\|\le R_*\) hold, the four updates take the linear form
\[
\bu^{k+1}-\bu_c=\mathbb T(\bu^k-\bu_c),
\] with $\mathbb T$ defined as 
\begin{equation}
\label{eq:rot-matrix}
\mathbb T=\begin{pmatrix}
1+\eta a-\tfrac{\eta}{\gamma} & -\eta b & \tfrac{\eta}{\gamma} & 0\\[2pt]
\tau b\bigl(1+\eta a-\tfrac{\eta}{\gamma}\bigr) & 1-\tfrac{\tau}{s}-\eta\tau b^2
  & \tfrac{\eta\tau b}{\gamma} & \tfrac{\tau}{s}\\[2pt]
\alpha\bigl(1+\eta a-\tfrac{\eta}{\gamma}\bigr) & -\alpha\eta b
  & 1-\alpha+\tfrac{\alpha\eta}{\gamma} & 0\\[2pt]
\beta\tau b\bigl(1+\eta a-\tfrac{\eta}{\gamma}\bigr)
  & \beta\bigl(1-\tfrac{\tau}{s}-\eta\tau b^2\bigr)
  & \tfrac{\beta\eta\tau b}{\gamma} & 1-\beta+\tfrac{\beta\tau}{s}
\end{pmatrix}.
\end{equation}
Under the conditions above, analyzing the LRALM iterates reduces to
studying the spectral properties of \(\mathbb T\).
The core idea is to construct a bounded trajectory whose KKT residual remains
above \(\epsilon\). We choose \(a\) and \(b\) so that \(\mathbb T\) has a
non-real eigenvalue of modulus one, which gives a component that does not
decay.

The construction in~\cite[Section~5.1.1]{li2026smoothing} chooses both the
hard instance and the initial point, with the latter lying along an
eigenvector of the iteration matrix. In our setting, the initial point is
fixed at zero, and we choose the instance alone. We therefore need to
obtain the required trajectory from this prescribed initial point. To this
end, we choose the centers \(x_c\) and \(\lambda_c\) in the instance so that
the shifted initial state is the sum of an eigenvector of \(\mathbb T\)
and a decaying component. The following lemma gives the spectral
properties and this decomposition for the parameters defined
in Appendix~\ref{app:lb-rotation-params}.

\begin{lemma}
\label{lem:rot-spectral}
Under the conditions of Lemma~\ref{lem:rot-estimates}, let
\(\mathbb T\) be the matrix in~\eqref{eq:rot-matrix}. Then:
\begin{enumerate}[label=(\roman*)]
\item \(\mathbb T\bv=\zeta\bv\), where \(|\zeta|=1\) and \(\zeta\notin\RR\);
\item \(\br\) lies in a \(\mathbb T\)-invariant subspace on which the spectral
  radius of \(\mathbb T\) is less than one;
\item \(\bv+\br=-\bu_c\).
\end{enumerate}
\end{lemma}
 
By Lemma~\ref{lem:rot-spectral}, the linear recursion initialized at zero
has the explicit solution
\[
\bu^k-\bu_c=\zeta^k\bv+\mathbb T^k\br.
\]
We next bound \(\mathbb T^k\br\) uniformly in \(k\) to verify that this
sequence stays in the quadratic region, both proximal outputs vanish,
and the KKT residual remains above \(\epsilon\).

\begin{lemma}
\label{lem:rot-estimates}
Let \(\alpha>\alpha_\epsilon\), \(0<\beta\le1\), \(0<\epsilon<c_2\), and
\(\epsilon(\rho+s)\le c_1\), with \(c_1\), \(c_2\), and \(\alpha_\epsilon\)
given in~\eqref{eq:rot-thresholds}. Then:
\begin{enumerate}[label=(\roman*)]
\item the instance of Example~\ref{ex:rotation} belongs to
  \(\mathcal I(\mathbf C_0,\Delta)\);
\item for all \(k\ge0\),
\[
\|(\mathbb T^k\br)_x\|\le\frac X4,
\qquad
\|(\mathbb T^k\br)_\lambda\|\le
s\Bigl(\frac{\gamma}{\eta}+\frac54\Bigr)bX.
\]
\end{enumerate}
\end{lemma}
The proofs of Lemmas~\ref{lem:rot-spectral} and~\ref{lem:rot-estimates} are
given in Appendix~\ref{app:lb-rotation}.

Using these estimates, we now show that the sequence generated by the linear
recursion coincides with the LRALM iterates and that its KKT residual
remains above \(\epsilon\) at every iteration.

\begin{lemma}
\label{lem:rot-trajectory}
Under the conditions of Lemma~\ref{lem:rot-estimates}, the iterates of
\(\mathrm{LRALM}_{\theta,\alpha,\beta}\) on Example~\ref{ex:rotation}
satisfy, for all \(k\ge0\),
\[
x^k-x_c=\zeta^kX+(\mathbb T^k\br)_x,
\qquad
\lambda^k-\lambda_c=\zeta^k\Lambda+(\mathbb T^k\br)_\lambda,
\]
and \(R_k\ge\frac76\epsilon\). Hence
\(T_\epsilon(\theta,\alpha,\beta;I,\zero)=+\infty\).
\end{lemma}

\begin{proof}
We prove by induction that, for all \(k\ge0\),
\[
\bu^k-\bu_c=\zeta^k\bv+\mathbb T^k\br.
\]
\emph{Base case \(k=0\).} By Lemma~\ref{lem:rot-spectral}~(iii),
\(\bv+\br=-\bu_c\). Since \(\bu^0=\zero\), the identity holds for \(k=0\).

\emph{Inductive step.} Assume that the identity holds for some \(k\ge0\).
Since \(|\zeta|=1\), Lemma~\ref{lem:rot-estimates}~(ii) gives
\[
\frac34X\le\|x^k-x_c\|\le\frac54X=R_*.
\]
Thus, \(x^k\) lies in the quadratic region. Moreover, the induction
hypothesis and Lemma~\ref{lem:rot-spectral}~(iii) imply
\(\|\lambda^k-\lambda_c\|\le\|\Lambda\|+\|(\mathbb T^k\br)_\lambda\|\)
and \(\|\lambda_c\|\le\|\Lambda\|+\|\br_\lambda\|\).
Lemma~\ref{lem:rot-estimates}~(ii) and the parameter bounds in
Lemma~\ref{lem:rot-parameter-bounds} therefore yield
\[
\|\lambda^k+\rho\calA(x^k)\|
\le2\|\Lambda\|+2s\Bigl(\frac{\gamma}{\eta}+\frac54\Bigr)bX+\frac54\rho bX
\le\frac13\Bigl(2+\frac14\Bigr)=\frac34.
\]
Hence \((x^k,\lambda^k)\in\mathcal R_\rho\). By
Proposition~\ref{prop:reduction}, the primal update is given by the first
row of~\eqref{eq:rot-matrix}. Using the induction hypothesis and
\(\mathbb T\bv=\zeta\bv\), we obtain
\[
x^{k+1}-x_c=\zeta^{k+1}X+(\mathbb T^{k+1}\br)_x.
\]
Lemma~\ref{lem:rot-estimates}~(ii) then gives
\(\frac34X\le\|x^{k+1}-x_c\|\le\frac54X\).
The same estimate as above, with \(x^{k+1}\) in place of \(x^k\), yields
\[
\|\lambda^k+\rho\calA(x^{k+1})\|\le\frac34.
\]
Thus, \((x^{k+1},\lambda^k)\in\mathcal R_\rho\) as well, so
Proposition~\ref{prop:reduction} also applies to the dual update.
Together with the center updates, this gives
\[
\bu^{k+1}-\bu_c
=\mathbb T(\bu^k-\bu_c)
=\zeta^{k+1}\bv+\mathbb T^{k+1}\br.
\]
This completes the induction. In particular, both proximal outputs vanish
at every iteration, and all primal iterates lie in the prescribed annulus.

Finally, we bound the residual. By Proposition~\ref{prop:embedding},
\[
R_k=\max\bigl\{\|{-a}(x^{k+1}-x_c)+b(\lambda^k-\lambda_c)\|,\
b\|x^{k+1}-x_c\|\bigr\}.
\]
The second term is at least \(\frac34bX\). For the first term,
Lemma~\ref{lem:rot-estimates}~(ii) and Lemma~\ref{lem:rot-parameter-bounds} give
\[
\|{-a}(x^{k+1}-x_c)+b(\lambda^k-\lambda_c)\|
\ge b\|\Lambda\|-s\Bigl(\frac{\gamma}{\eta}+\frac54\Bigr)b^2X-\frac54aX
\ge\frac{27}{32}b\|\Lambda\|-s\Bigl(\frac{\gamma}{\eta}+\frac52\Bigr)b^2X .
\]
If \(\|\Lambda\|\ge X\), Lemma~\ref{lem:rot-parameter-bounds} bounds the
right-hand side below by \(\frac{23}{32}b\|\Lambda\|\). If
\(X\ge\|\Lambda\|\), we use the second term instead. In both cases,
\(R_k\ge\frac b2\max\{X,\|\Lambda\|\}=\frac b6\ge\frac76\epsilon\).
\end{proof}

\subsection{Proof of Lemma~\ref{lem:main-huber}}
\label{subsec:lb-huber}

\begin{proof}[Proof of Lemma~\ref{lem:main-huber}]
Since \(h\) is convex and \(1\)-Lipschitz, \(\calA\) is constant, and
\(F=f+h(a_0)\) has the same initial gap \(\Delta\) as \(f\), we have
\((f,\calA,h)\in\mathcal I(\mathbf C_0,\Delta)\).
Run \(\mathrm{LRALM}_{\theta,\alpha,0}\) on Example~\ref{ex:huber} from
\(\bu^0=\zero\), and set
\[
\varphi:=1-\frac{\tau}{s}-\frac{\tau}{\rho+q_h},
\qquad
\bar\lambda:=\frac{s\,q_ha_0}{\rho+q_h+s}.
\]
We first collect three bounds on the parameters. Since \(0<\tau\le s\) and
\(\rho+q_h>2s\),
\begin{equation}
\label{eq:huber-ratio}
-\frac12<-\frac{\tau}{\rho+q_h}\le\varphi<1,
\qquad
\frac{s}{\rho+q_h}<\frac12 .
\end{equation}
Hence \(\varphi^k\ge-\frac12\) and \(0\le1-\varphi^k\le2\) for all
\(k\ge0\). Since \(2s\le\rho+q_h+s\), it follows that
\(0\le\bar\lambda(1-\varphi^k)\le2\bar\lambda\le q_ha_0\). Finally,
\(s\ge\tau\) gives \(1\le(\rho+s)/\tau\), so the third term in the
definition of \(c_1\) yields
\[
q_ha_0=4\epsilon(\rho+3s+1)
\le4\epsilon\bigl(3+\tau^{-1}\bigr)(\rho+s)
\le4c_1\bigl(3+\tau^{-1}\bigr)
\le\frac14 .
\]

We first show by induction that for all \(k\ge0\),
\(\mu^k=0\), \(\lambda^k=\bar\lambda(1-\varphi^k)\), and
\(y^{k+1}=(\rho a_0+\lambda^k)/(\rho+q_h)\). Since \(\beta=0\) and \(\mu^0=0\),
the dual center update gives \(\mu^k=0\) for all \(k\). Next, fix \(x\in\RR\)
and \(\lambda\in[0,q_ha_0]\), and set
\(y_\lambda:=(\rho a_0+\lambda)/(\rho+q_h)\). Then
\(0\le y_\lambda\le a_0\le q_h^{-1}\), where the last inequality uses
\(q_ha_0\le\frac14\). On \([-q_h^{-1},q_h^{-1}]\), the function \(h\) is
differentiable with \(\nabla h(y)=q_hy\), and \(y_\lambda\) satisfies
\[
q_hy_\lambda-\lambda-\rho\bigl(a_0-y_\lambda\bigr)=0,
\]
which is the optimality condition of the strongly convex problem
in~\eqref{eq:reduced-al}. Hence \(y(x,\lambda)=y_\lambda\), and
\[
\calA(x)-y(x,\lambda)=\frac{q_ha_0-\lambda}{\rho+q_h}.
\]
By Danskin's theorem,
\(\nabla_\lambda\calD_\rho(x,\lambda)=\calA(x)-y(x,\lambda)\).

The formula for \(\lambda^k\) holds at \(k=0\) because \(\lambda^0=0\).
Assume that it holds for some \(k\ge0\). Then \(\lambda^k\in[0,q_ha_0]\), so
\(y^{k+1}=y(x^{k+1},\lambda^k)=y_{\lambda^k}\). Together with \(\mu^k=0\), the
dual update in~\eqref{eq:splalm-updates} reads
\[
\lambda^{k+1}
=\lambda^k+\tau\Bigl[\frac{q_ha_0-\lambda^k}{\rho+q_h}
-\frac{\lambda^k}{s}\Bigr]
=\varphi\lambda^k+\frac{\tau q_ha_0}{\rho+q_h}.
\]
A direct computation shows that
\(\bar\lambda=\varphi\bar\lambda+\tau q_ha_0/(\rho+q_h)\). Subtracting the
two identities gives
\(\lambda^{k+1}-\bar\lambda=\varphi(\lambda^k-\bar\lambda)
=-\bar\lambda\varphi^{k+1}\), completing the induction.

Finally, since \(J_{\calA}=0\), \eqref{eq:kkt-residual} gives
\(R_k=\max\{\|\nabla f(x^{k+1})\|,\|\calA(x^{k+1})-y^{k+1}\|\}\). Using
\(q_ha_0-\bar\lambda=q_ha_0(\rho+q_h)/(\rho+q_h+s)\),
\[
\calA(x^{k+1})-y^{k+1}
=\frac{q_ha_0-\lambda^k}{\rho+q_h}
=\frac{q_ha_0}{\rho+q_h+s}\Bigl(1+\frac{s}{\rho+q_h}\varphi^k\Bigr)
=4\epsilon\Bigl(1+\frac{s}{\rho+q_h}\varphi^k\Bigr),
\]
where the last equality uses \(\rho+q_h+s=\rho+3s+1\).
By~\eqref{eq:huber-ratio}, \(\frac{s}{\rho+q_h}\varphi^k>-\frac14\), so
\(R_k\ge\|\calA(x^{k+1})-y^{k+1}\|>3\epsilon\). Thus no output is an
\(\epsilon\)-KKT point, and
\(T_\epsilon(\theta,\alpha,0;I,\zero)=+\infty\).
\end{proof}

\subsection{Details for the rotating instance}
\label{app:lb-rotation}

\subsubsection{Parameters}
\label{app:lb-rotation-params}
Let \(V:=\beta+\frac{\tau}{s}(1-\beta)\) and
\[
\nu:=\frac{64\eta\tau\alpha\beta\epsilon^2}
  {\bigl(\alpha+\frac{\eta}{\gamma}(1-\alpha)\bigr)V},
\qquad
\zeta:=1-\frac{\nu}{2}+i\sqrt{\nu-\frac{\nu^2}{4}} .
\]
Under the conditions of Lemma~\ref{lem:rot-estimates}, \(0<\nu<4\); hence
\(\zeta\notin\RR\), \(|\zeta|=1\), and \(|\zeta-1|^2=\nu\). Let
\(W:=1+\frac{\eta}{\gamma}\alpha(1-\alpha)/|\zeta-1+\alpha|^2\), and define
\[
b^2:=\frac{\nu W|\zeta-1+V|^2}{\eta\tau\bigl(\beta V+(1-V)\nu\bigr)},
\qquad
a:=\frac{\nu(W-1)}{\eta\alpha}+\frac{\tau b^2(V-\beta)}{|\zeta-1+V|^2},
\qquad
c_\kappa:=\rho b^2+a.
\]
Let \(X>0\) be such that \(\max\{X,\|\Lambda\|\}=\frac13\), where
\(\Lambda:=\frac{\tau b\zeta(\zeta-1+\beta)}{(\zeta-1)(\zeta-1+V)}X\), and set
\begin{equation}
\bv:=\Bigl(X,\ \Lambda,\ \frac{\alpha\zeta}{\zeta-1+\alpha}X,\
  \frac{\beta\zeta}{\zeta-1+\beta}\Lambda\Bigr).
\label{eq:rot-eigvec}
\end{equation}
The shifts \(x_c\) and \(\lambda_c\) are given in~\eqref{eq:rot-shift}. They
are chosen so that \(-\bu_c-\bv\) lies in the \(\mathbb T\)-invariant subspace
on which the spectral radius of \(\mathbb T\) is less than one.


We use the following choices of the constants in Theorem~\ref{thm:lower-bound}
and of the partition threshold \(\alpha_\epsilon\):
\begin{equation}
\label{eq:rot-thresholds}
\begin{aligned}
c_1&:=\min\left\{
\frac1{32}\sqrt{\frac{\tau}{\eta(1+4\gamma/\eta)}},\,
\frac1{64(\gamma/\eta+5/2)},\,
\frac1{16(3+\tau^{-1})}\right\},\\
c_2&:=\min\left\{\frac1{64},\,
\frac1{32\sqrt{\max\{1,4\tau,4\gamma\tau\}}},\,
\frac\Delta{12},\,\frac{\sqrt\Delta}{2},\,
\sqrt{\frac\Delta{64\gamma}}\right\},\\
\alpha_\epsilon&:=1024\max\{1,4\tau,4\gamma\tau\}\,\epsilon^2.
\end{aligned}
\end{equation}
Throughout this appendix, assume the parameter conditions of
Lemma~\ref{lem:rot-estimates}. All complex scalars represent vectors in
\(\RR^2\), with the Euclidean norm equal to the complex modulus.
For the calculations below, write
\[
\chi:=\frac\eta\gamma,\qquad d_\tau:=\frac\tau s,\qquad
U:=\alpha+\chi(1-\alpha),\qquad D:=\alpha W,
\]
\[
E:=|\zeta-1+\alpha|^2=\alpha^2+(1-\alpha)\nu,\qquad
H:=|\zeta-1+V|^2=V^2+(1-V)\nu,
\]
\[
\Upsilon:=\beta V+(1-V)\nu,\qquad
\xi:=\frac{\eta\tau b^2(V-\beta)}H,\qquad
C_\mathrm{tr}:=\frac1{d_\tau}+\frac{2(1+d_\tau^{-1})}{\chi}.
\]
These quantities are used only to prove the estimates and spectral identities.

\subsubsection{Parameter bounds and proof of Lemma~\ref{lem:rot-estimates}~(i)}
\label{app:lb-rotation-params-proof}
\begin{lemma}
\label{lem:rot-parameter-bounds}
Under the conditions of Lemma~\ref{lem:rot-estimates}, the parameters satisfy
\[
7\epsilon\le b\le8\epsilon,
\qquad \max\{X,\|\Lambda\|\}=\frac13,
\]
\[
aX\le\frac18b\|\Lambda\|+sb^2X,
\]
\[
bs\Bigl(\frac{\gamma}{\eta}+\frac52\Bigr)\le\frac18,
\qquad
b\Bigl(2s\Bigl(\frac{\gamma}{\eta}+\frac54\Bigr)+\frac54\rho\Bigr)\le\frac14.
\]
\end{lemma}

\begin{proof}[Proof of Lemma~\ref{lem:rot-parameter-bounds}]
We first establish the bounds needed in all three proofs. Since
\(U\ge\chi\), \(V\ge\beta\), and
\(\alpha>1024\max\{1,4\tau,4\gamma\tau\}\epsilon^2\),
\[
0<\frac\nu{\alpha^2}
\le\frac{64\gamma\tau\epsilon^2}{\alpha}<\frac1{64}.
\]
In particular, \(0<\nu<4\), so \(\zeta\) is well defined and
\(|\zeta|=1\). The definition of \(W\) gives
\[
\frac{64}{65}U\le
D=\alpha+\frac{\chi(1-\alpha)}{1+(1-\alpha)\nu/\alpha^2}
\le U\le1,
\qquad D\ge\frac\chi2.
\]
Since \(V\Upsilon-\beta H=(1-V)(V-\beta)\nu\ge0\),
\[
\frac{b^2}{64\epsilon^2}=\frac D U\frac{\beta H}{V\Upsilon}\le1.
\]
Also, \(H\ge V^2\), \(V\ge d_\tau\), and \(V-\beta\le V\) imply
\(0\le\xi\le\eta s b^2\). Using
\(C_\mathrm{tr}\le(1+4/\chi)/d_\tau\) and~\eqref{eq:rot-thresholds}, we obtain
\begin{equation}
\label{eq:rot-small-feedback}
\xi C_\mathrm{tr}
\le\frac{64\eta(1+4/\chi)}\tau(\epsilon s)^2
\le\frac1{16}.
\end{equation}
Thus \(\xi/d_\tau\le1/16\). The identity
\[
\frac{\beta H}{V\Upsilon}
=1-\frac{(1-V)\xi}{WV}\ge\frac{15}{16}
\]
now yields \(b^2/(64\epsilon^2)\ge12/13>49/64\), proving
\(7\epsilon\le b\le8\epsilon\).
The normalization of \(X\) gives \(\max\{X,\|\Lambda\|\}=1/3\).

To bound \(a\), set
\(a_\mathrm{p}:=\nu(W-1)/(\eta\alpha)\) and \(a_\mathrm{d}:=\xi/\eta\).
The definitions imply
\begin{equation}
\label{eq:rot-algebra-identities}
\eta a=\frac{\nu(W-1)}\alpha+\xi,
\qquad \eta\tau b^2=\nu W+V\xi,
\qquad U-D=\frac{(1-\alpha)\nu(W-1)}\alpha.
\end{equation}
Hence \(a=a_\mathrm{p}+a_\mathrm{d}\),
\(a_\mathrm{p}\le\tau b^2/\alpha\), and \(a_\mathrm{d}\le sb^2\).
Since \(c_1\le1/224\) and \(\epsilon<1/64\),
\[
0<c_\kappa\le(\rho+s)b^2+\frac{\tau b^2}\alpha
\le64c_1\epsilon+\frac{\tau}{16\max\{1,4\tau,4\gamma\tau\}}
\le\frac1{32}.
\]
For the amplitude bounds, put \(E_\beta:=\beta^2+(1-\beta)\nu\).
The formula for \(\Lambda\) gives
\(\|\Lambda\|^2/X^2=\tau WE_\beta/(\eta\Upsilon)\).
Moreover,
\[
\Upsilon\le
\operatorname{Re}\bigl((\zeta-1+V)\overline{(\zeta-1+\beta)}\bigr)
\le\sqrt{HE_\beta},
\]
where the difference in the first inequality is \((V-\beta)\nu/2\).
Consequently,
\[
\frac{a_\mathrm{p}X}{b\|\Lambda\|}
=\frac{\sqrt\nu(W-1)\Upsilon}{\alpha W\sqrt{HE_\beta}}
\le\frac{\sqrt\nu}\alpha\le\frac18,
\]
\[
\frac{a_\mathrm{p}X^2}{b^2\|\Lambda\|^2}
=\frac{\eta(W-1)\Upsilon^2}{\alpha W^2HE_\beta}
\le\frac\eta D\le2\gamma.
\]
Adding the bound on \(a_\mathrm{d}\) proves the asserted bound on \(aX\)
and also gives
\[
c_\kappa X^2\le(\rho+s+2\gamma)b^2\max\{X^2,\|\Lambda\|^2\}.
\]
Finally, \(b\le8\epsilon\) and~\eqref{eq:rot-thresholds} imply
\[
bs\Bigl(\frac1\chi+\frac52\Bigr)
\le8c_1\Bigl(\frac1\chi+\frac52\Bigr)\le\frac18,
\]
\[
b\Bigl(2s\Bigl(\frac1\chi+\frac54\Bigr)+\frac54\rho\Bigr)
\le16c_1\Bigl(\frac1\chi+\frac54\Bigr)
\le\frac{1/\chi+5/4}{4(1/\chi+5/2)}<\frac14.
\]
\end{proof}

\begin{proof}[Proof of Lemma~\ref{lem:rot-estimates}~(i)]
We use the bounds on \(\mathbb T^k\br\) proved in
Appendix~\ref{app:lb-rotation-transient-proof}; that proof uses only
Lemma~\ref{lem:rot-parameter-bounds} and the spectral identities, and does
not use membership in the problem class. At \(k=0\), those bounds and
\eqref{eq:rot-shift} give
\[
\|x_c\|\le\frac54X=R_*,\qquad
\|\lambda_c\|\le\|\Lambda\|+s\Bigl(\frac1\chi+\frac54\Bigr)bX
\le\frac13\left(1+\frac{1/\chi+5/4}{8(1/\chi+5/2)}\right)<\frac12.
\]
The function \(\psi\) and its first derivative agree at \(R_*\) and
\(3R_*\). On each smooth radial piece, the Hessian of
\(\psi(\|x-x_c\|)\) has radial eigenvalue \(\psi''(r)\) and tangential
eigenvalue \(\psi'(r)/r\), both of absolute value at most \(c_\kappa\).
Near \(x_c\), this function is quadratic. Integrating its continuous
gradient along line segments therefore shows that
\(\operatorname{Lip}(\nabla f)\le c_\kappa\le1/32\).
Also \(|\psi'(r)|\le c_\kappa R_*\) for every \(r\ge0\), so
\[
\|\nabla f(x)\|\le c_\kappa R_*+b\|\lambda_c\|
\le\frac5{384}+\frac12<1.
\]
Here we used \(R_*\le5/12\) and \(b\le8\epsilon<1\).
The affine mapping \(\calA\) has constant Jacobian of norm \(b\le1\),
and \(h=\|\cdot\|_2\) is convex and 1-Lipschitz.

It remains to bound the initial gap. Since \(\psi(r)\ge-c_\kappa R_*^2\)
and \(\|\lambda_c\|<1/2\),
\[
F(x)=\psi(\|x-x_c\|)-b\langle\lambda_c,x-x_c\rangle+b\|x-x_c\|
\ge-c_\kappa R_*^2.
\]
At the origin, \(\|x_c\|\le R_*\) and \(\psi(\|x_c\|)\le0\), whence
\[
F(0)-\inf F\le\frac32bR_*+c_\kappa R_*^2
\le5\epsilon+\frac{100}{9}(\rho+s+2\gamma)\epsilon^2.
\]
By \(\epsilon(\rho+s)\le c_1\le1/224\), this is at most
\(6\epsilon+32\gamma\epsilon^2\le\Delta\), where the last inequality
uses \(\epsilon<c_2\). Thus the instance belongs to
\(\mathcal I(\mathbf C_0,\Delta)\).
\end{proof}

\subsubsection{Proof of Lemma~\ref{lem:rot-spectral}}
\label{app:lb-rotation-spectral-proof}
Define
\[
\begin{aligned}
r_x&:=-\frac{(W-1)\frac{\zeta-1}{\zeta-1+\alpha}+\frac{\xi\zeta}{\zeta-1+V}}
  {W+\xi}\,X,\\
r_\lambda&:=\frac{1}{\eta b}\Bigl[\frac{\eta}{\gamma}
  \frac{(1-\alpha)(\zeta-1)}{\zeta-1+\alpha}X+(\eta a+\alpha W)r_x
  -\xi\Bigl((1-\alpha)r_x+\frac{(V-\alpha)\zeta}{\zeta-1+V}X\Bigr)\Bigr].
\end{aligned}
\]
Set
\begin{equation}
x_c:=-(X+r_x),\qquad \lambda_c:=-(\Lambda+r_\lambda),\qquad
\br:=-\bu_c-\bv .
\label{eq:rot-shift}
\end{equation}
\begin{proof}
For this calculation, write \(Z:=\bv_z\), \(M:=\bv_\mu\), and
\(\vartheta:=(1-V)\xi\). The characteristic polynomial factors as
\begin{equation}
\label{eq:rot-characteristic-factorization}
\begin{aligned}
\det((1+t)I-\mathbb T)
&=t^2(t+U)(t+V)-\eta a\,t(t+\alpha)(t+V)\\
&\quad+\eta\tau b^2(t+1)(t+\alpha)(t+\beta)\\
&=(t^2+\nu t+\nu)
  \bigl((t+D)(t+V)-\vartheta(t+\alpha)\bigr).
\end{aligned}
\end{equation}
The first equality follows by eliminating the two center coordinates in
\eqref{eq:rot-matrix}; expansion using~\eqref{eq:rot-algebra-identities}
gives the second equality. Since
\((\zeta-1)^2+\nu(\zeta-1)+\nu=0\), substitution of
\eqref{eq:rot-eigvec} in the four rows of \(\mathbb T\) gives
\[
\mathbb T\bv=\zeta\bv.
\]
For example, the center rows follow from
\((\zeta-1+\alpha)Z=\alpha\zeta X\) and
\((\zeta-1+\beta)M=\beta\zeta\Lambda\); the dual row reduces to
\((\zeta-1)(\zeta-1+V)\Lambda
=\tau b\zeta(\zeta-1+\beta)X\).
The primal row then follows from the first factor in
\eqref{eq:rot-characteristic-factorization}. The bounds on \(\nu\)
proved above give \(\zeta\notin\RR\) and \(|\zeta|=1\).

We next prove
\begin{equation}
\label{eq:rot-annihilator}
Q_\mathrm{f}(\mathbb T)\br=0,
\qquad Q_\mathrm{f}(t):=(t-1+D)(t-1+V)-\vartheta(t-1+\alpha).
\end{equation}
First consider distinct roots \(t_1,t_2\) of \(Q_\mathrm{f}\), with
\(\alpha<1\), \(\beta<1\), \(d_\tau<1\), \(V\ne\alpha\), and all
following denominators nonzero. An eigenvector associated with a root \(t\) is
\[
\left(1,\ \frac{\tau bt(t-1+\beta)}{(t-1)(t-1+V)},\
\frac{\alpha t}{t-1+\alpha},\
\frac{\tau b\beta t^2}{(t-1)(t-1+V)}\right)^\top.
\]
Let its coefficients at \(t_1,t_2\) be \(c_1',c_2'\).
Requiring that their sum with \(\bv\) have equal first and third
coordinates, and equal second and fourth coordinates, gives
\[
\sum_{j=1}^2c_j'\frac{t_j-1}{t_j-1+\alpha}
=-\frac{\zeta-1}{\zeta-1+\alpha}X,
\qquad
\sum_{j=1}^2c_j'\frac{t_j}{t_j-1+V}
=-\frac\zeta{\zeta-1+V}X.
\]
To solve these two equations, set
\(S_0:=c_1'+c_2'\), \(S_1:=\sum_j(t_j-1)c_j'\),
\(C_\mathrm{f}:=D+V-\vartheta\), and \(D_\mathrm{f}:=DV-\alpha\vartheta\).
At a root \(v=t_j-1\) of \(v^2+C_\mathrm{f}v+D_\mathrm{f}\),
\[
\frac v{v+\alpha}
=\frac{D_\mathrm{f}+\alpha v}{\alpha^2-C_\mathrm{f}\alpha+D_\mathrm{f}},
\qquad
\frac{v+1}{v+V}
=\frac{D_\mathrm{f}-C_\mathrm{f}+V-(1-V)v}{V^2-C_\mathrm{f}V+D_\mathrm{f}}.
\]
The two denominators equal \(\alpha(W-1)(V-\alpha)\) and
\(\vartheta(V-\alpha)\), respectively. The resulting system has
coefficient determinant \(\alpha(1-V)(\alpha-V)(W+\xi)\ne0\), and its solution is
\[
S_0=r_x,\qquad
S_1=-Dr_x+\xi\left((1-\alpha)r_x+
\frac{(V-\alpha)\zeta}{\zeta-1+V}X\right).
\]
The primal row of \(\mathbb T\) now gives
\(\eta br_\lambda=\chi(X-Z)+\eta ar_x-S_1\), which is the stated
formula for \(r_\lambda\). The center synchronization gives
\(r_z=X-Z+r_x\) and \(r_\mu=\Lambda-M+r_\lambda\).
Thus this linear combination is precisely \(\br\), proving
\eqref{eq:rot-annihilator} in the stated nondegenerate case.

All entries of \(Q_\mathrm{f}(\mathbb T)\br\) are continuous functions
of the admissible parameters: \(\alpha,b,V>0\), \(W+\xi>0\), and
\(|\zeta-1+\alpha|,|\zeta-1+V|>0\).
To justify the density used here, first take
\(\alpha,\beta,d_\tau\in(0,1)\) and \(V\ne\alpha\).
Then \(W-1>0\) and \(\vartheta>0\), so the two displayed denominators
and the coefficient determinant are nonzero. Moreover,
\[
Q_\mathrm{f}(1)=DV-\alpha\vartheta
=\frac{\alpha\beta\eta\tau b^2}{\nu}>0,
\]
so neither root equals one. The remaining exclusion is a repeated root.
Given \(\alpha,\beta,d_\tau\), its discriminant
\[
(D+V-\vartheta)^2-4(DV-\alpha\vartheta)
\]
is a rational function of \(\epsilon^2\) and tends to \((U-V)^2\) as
\(\epsilon\downarrow0\). If \(U\ne V\), this rational function is not
identically zero, so its zeros are isolated. Since
\(V=\beta+d_\tau(1-\beta)\) varies strictly with \(\beta\) in the interior,
the restrictions \(V\ne\alpha\) and \(V\ne U\) also define a dense subset.
These perturbations can be made within the admissible parameter set:
first decrease \(\epsilon\) slightly if
\(\epsilon(\rho+s)=c_1\), and then perturb the other parameters.
Taking limits therefore proves~\eqref{eq:rot-annihilator} also when roots
coincide, \(V=\alpha\), or \(\alpha=1\), \(\beta=1\), or \(d_\tau=1\).

To locate the roots of \(Q_\mathrm{f}\), note that
\(0<D,V\le1\) and, by~\eqref{eq:rot-small-feedback},
\[
\frac{2\vartheta}{DV}\le\frac{4\xi}{\chi d_\tau}
\le2\xi C_\mathrm{tr}\le\frac18.
\]
If \(|t|\ge1\), then
\[
|(t-1+D)(t-1+V)|\ge DV|t|^2
>2\vartheta|t|\ge|\vartheta(t-1+\alpha)|.
\]
Thus every root of \(Q_\mathrm{f}\) lies inside the unit circle.
The subspace \(\ker Q_\mathrm{f}(\mathbb T)\) is invariant under \(\mathbb T\),
and every eigenvalue of the restriction is a root of \(Q_\mathrm{f}\).
It contains \(\br\) by~\eqref{eq:rot-annihilator}, proving (ii), including
repeated roots. Finally, \(\bv+\br=-\bu_c\) follows from
\eqref{eq:rot-shift}, proving (iii).
\end{proof}

\subsubsection{Proof of Lemma~\ref{lem:rot-estimates}~(ii)}
\label{app:lb-rotation-transient-proof}
\begin{proof}
Write \((\mathbb T^k\br)_x\), \((\mathbb T^k\br)_\lambda\), and
\((\mathbb T^k\br)_\mu\) as \(p_k,\ell_k,m_k\), respectively.
These are components of the linear recursion, not new algorithm iterates.
By the preceding spectral proof, all these components tend to zero.
Put \(r_\mathrm{p}:=1-D\), \(r_\mathrm{d}:=1-V\), and
\(\vartheta:=(1-V)\xi\). Equation~\eqref{eq:rot-annihilator} gives
\[
p_{k+2}-(r_\mathrm{p}+r_\mathrm{d})p_{k+1}+r_\mathrm{p}r_\mathrm{d}p_k
=\vartheta\bigl(p_{k+1}-(1-\alpha)p_k\bigr).
\]
The initial values are \(p_0=r_x\) and
\(p_1=r_\mathrm{p}r_x+\omega_x\), where
\[
\omega_x:=\xi\left((1-\alpha)r_x+
\frac{(V-\alpha)\zeta}{\zeta-1+V}X\right).
\]
Since \(|\zeta-1|=\sqrt\nu\), the formula for \(r_x\) implies
\[
|r_x|\le\left(\frac{\sqrt\nu}{\alpha}+\frac\xi{d_\tau}\right)X
\le\frac{3X}{16}.
\]
As \(|V-\alpha|\le1\), it follows that
\(|\omega_x|\le\xi(1+d_\tau^{-1})X\). Hence
\begin{equation}
\label{eq:rot-initial-transient}
|r_x|+\frac{|\omega_x|}D
\le\left(\frac18+\xi C_\mathrm{tr}\right)X\le\frac{3X}{16}.
\end{equation}

For \(j\ge0\), define the convolution kernel
\(G_j:=\sum_{i=0}^jr_\mathrm{p}^{j-i}r_\mathrm{d}^i\).
Since \(0\le r_\mathrm{p},r_\mathrm{d}<1\),
\[
0\le G_j\le\sum_{i=0}^jr_\mathrm{p}^{j-i}\le\frac1D,
\qquad \sum_{j\ge0}G_j=\frac1{DV}.
\]
Solving the scalar recurrence gives, for \(k\ge1\),
\[
p_k=r_\mathrm{p}^kr_x+G_{k-1}\omega_x+
\vartheta\sum_{j=0}^{k-2}G_{k-2-j}\bigl(p_{j+1}-(1-\alpha)p_j\bigr),
\]
with an empty sum when \(k=1\). This formula follows by substitution
in the recurrence and its two initial values.
For each finite \(N\), let \(M_N:=\max_{0\le k\le N}|p_k|\).
Using~\eqref{eq:rot-initial-transient} and \(2\vartheta/(DV)\le1/4\),
\[
M_N\le\frac{3X}{16}+\frac{2\vartheta}{DV}M_N,
\qquad M_N\le\frac X4.
\]
Since \(N\) is arbitrary, this proves the primal bound.
Summing absolute values on a finite prefix instead gives
\[
\sum_{k=0}^N|p_k|
\le\frac{|r_x|}D+\frac{|\omega_x|}{DV}
+\frac{2\vartheta}{DV}\sum_{k=0}^N|p_k|.
\]
Because \(D,V\le1\),
\[
\frac{|r_x|}D+\frac{|\omega_x|}{DV}
\le\frac{|r_x|+|\omega_x|/D}{DV}.
\]
Letting \(N\) increase yields
\begin{equation}
\label{eq:rot-transient-sum}
\sum_{k\ge0}|p_k|\le\frac X{4DV}\le\frac X{\chi d_\tau}.
\end{equation}

For the dual bound, introduce the difference and weighted average
\[
e_k:=\ell_k-m_k,
\qquad q_k:=\frac{\beta(1-d_\tau)\ell_k+d_\tau m_k}{V}.
\]
The dual and center rows of \(\mathbb T\) give
\[
e_{k+1}=(1-V)e_k+\tau b(1-\beta)p_{k+1},
\qquad q_{k+1}=q_k+\frac{\tau b\beta}{V}p_{k+1}.
\]
The synchronization at \(k=0\) implies
\[
e_0=-\frac{\tau b(1-\beta)\zeta}{\zeta-1+V}X,
\qquad |e_0|\le sbX.
\]
Iterating the first recurrence and using \(|p_k|\le X/4\) gives
\[
|e_k|\le sbX+\frac{\tau b(1-\beta)}V\frac X4\le\frac54sbX.
\]
Since \(q_k\to0\), the second recurrence and the absolutely summable
bound~\eqref{eq:rot-transient-sum} give
\[
q_k=-\frac{\tau b\beta}V\sum_{j=k+1}^{\infty}p_j,
\qquad |q_k|\le\frac{s}{\chi}bX.
\]
Finally, \(\ell_k=q_k+(d_\tau/V)e_k\), with \(d_\tau/V\le1\).
Thus
\[
\|(\mathbb T^k\br)_\lambda\|=|\ell_k|
\le s\left(\frac1\chi+\frac54\right)bX
=s\left(\frac\gamma\eta+\frac54\right)bX,
\]
as required.
\end{proof}

\subsection{Proof of Lemma~\ref{lem:main-radial}}
\label{app:radial}

Throughout this subsection we identify \(\RR^2\) with \(\mathbb C\), so each
complex coordinate represents a vector in \(\RR^2\) with Euclidean norm equal to
its modulus and \(\langle p,q\rangle=\operatorname{Re}(\overline p q)\).
Given fixed \(\rho>0\), \(0<\eta\le\gamma\), \(\tau>0\), \(c_s>0\), and
\(\Delta>0\), set \(\chi:=\eta/\gamma\in(0,1]\) and define the constants
\begin{equation}
\begin{aligned}
\omega_0&:=\min\left\{
  \frac12,\;
  \frac{\sqrt{\eta\tau}}{2},\;
  \frac12\sqrt{\frac{\eta}{1+\rho/\tau}}\right\},
& A_0&:=\min\left\{\frac14,\;\frac{\Delta}{2},\;\frac{1}{2+\rho}\right\},\\
c_\mathrm{r}&:=A_0\omega_0
  \min\left\{\frac{1}{\sqrt{\eta\tau}},\;\frac{1}{2\tau}\right\},
& B&:=\max\left\{1,\;\frac{2\tau}{c_s\omega_0^2},\;\frac{1}{c_\mathrm{r}}\right\},\\
\epsilon_\mathrm{r}&:=\min\left\{\frac12,\;\frac{c_s}{2\tau},\;\frac{1}{B}\right\}.
\end{aligned}
\label{eq:radial-fixed-constants}
\end{equation}
These constants depend only on \((\rho,\gamma,\eta,\tau,c_s,\Delta)\).

Given \(\epsilon\in(0,\epsilon_\mathrm{r})\), \(\alpha\in(B\epsilon,1]\),
\(\beta\in[0,1]\), and \(s:=c_s/\epsilon\), and write \(d:=\tau/s=\tau\epsilon/c_s\).
By~\eqref{eq:radial-fixed-constants}, we have \(0<d\le\frac12\),
\(d<\frac{\omega_0^2}{2}\alpha\), and \(s>\tau\). Define the instance parameters
\begin{equation}
\begin{aligned}
\nu&:=\omega_0^2\alpha^2,
&\zeta&:=1-\frac{\nu}{2}+\mathrm{i}\sqrt{\nu-\frac{\nu^2}{4}},
&u&:=\zeta-1,\\
V&:=\beta+d(1-\beta),
&E_\alpha&:=|u+\alpha|^2=\alpha^2+(1-\alpha)\nu,
&W&:=1+\frac{\chi\alpha(1-\alpha)}{E_\alpha},\\
H&:=|u+V|^2=V^2+(1-V)\nu,
&\Upsilon&:=\beta V+(1-V)\nu,
&\xi&:=\frac{\nu W(V-\beta)}{\Upsilon},\\
b^2&:=\frac{\nu WH}{\eta\tau\Upsilon},
&a&:=\frac{\nu(W-1)}{\eta\alpha}+\frac{\xi}{\eta},
&c_\kappa&:=\rho b^2+a,
\end{aligned}
\label{eq:radial-parameters}
\end{equation}
with \(b>0\), and define the rotating eigenmode, centers, and initialization by
\begin{equation}
\begin{aligned}
L&:=\frac{\tau b\zeta(u+\beta)}{u(u+V)},
&X&:=\frac{A_0}{\max\{1,|L|\}},
&\Lambda&:=LX,\\
Z&:=\frac{\alpha\zeta}{u+\alpha}X,
&M&:=\frac{\beta\zeta}{u+\beta}\Lambda,
&x_c&:=-X,\quad\lambda_c:=-M,\\
\bv&:=(X,\Lambda,Z,M)^{\top},
&\bu_c&:=(x_c,\lambda_c,x_c,\lambda_c)^{\top},
&\bu^0&:=\bu_c+\bv=(0,\Lambda-M,Z-X,0)^{\top}.
\end{aligned}
\label{eq:radial-eigenmode}
\end{equation}

\begin{proof}[Proof of Lemma~\ref{lem:main-radial}]
We verify the parameter bounds and initialization admissibility, class
membership, exact trajectory, and residual lower bound in four steps.

\paragraph{Step 1: parameter bounds and \(\bu^0\in\mathcal S_0(\mathbf C_0)\).}
Since \(0<\nu\le\omega_0^2\le\frac14\), we have \(|\zeta|=1\),
\(|u|=\sqrt{\nu}=\omega_0\alpha>0\), and \(u^2+\nu u+\nu=0\).
For every \(t\in[0,1]\),
\[
|u+t|^2=t^2+(1-t)\nu=\Bigl(t-\frac{\nu}{2}\Bigr)^2+\nu\Bigl(1-\frac{\nu}{4}\Bigr)
\ge\frac34\nu>\frac14\nu>0.
\]
Hence \(E_\alpha\ge\alpha^2>0\), \(H\ge\frac34\nu>0\), and all denominators
in~\eqref{eq:radial-parameters}--\eqref{eq:radial-eigenmode} are nonzero
(for \(\beta=0\), \(V=d\le\frac12\) gives \(\Upsilon=(1-d)\nu>0\); for
\(\beta>0\), \(\Upsilon\ge\beta V>0\)).
Substituting \(\nu=\omega_0^2\alpha^2\) yields
\[
\alpha W=\alpha+\frac{\chi(1-\alpha)}{1+(1-\alpha)\omega_0^2}\le1,
\]
so \(W\ge1\) and \(\alpha(W-1)\le1\). If \(\beta=1\), then \(\xi=0\). If
\(\beta<1\), then \(V-\beta=d(1-\beta)\), \(1-V=(1-d)(1-\beta)\), and
\(\Upsilon\ge(1-V)\nu\) imply
\[
\xi\le\frac{d}{1-d}W\le2dW=\frac{2d}{\alpha}(\alpha W)\le\frac{2d}{\alpha}<\omega_0^2.
\]
Using \(H-\Upsilon=V(V-\beta)\), we obtain \(\eta\tau b^2=\nu W+V\xi\).
Since \(W\ge1\), \(\nu W=\omega_0^2\alpha(\alpha W)\le\omega_0^2\), and
\(V\le1\), it follows that
\[
\frac{\omega_0^2\alpha^2}{\eta\tau}\le b^2\le\frac{2\omega_0^2}{\eta\tau}\le\frac12,
\qquad
0\le\eta a=\omega_0^2\alpha(W-1)+\xi\le2\omega_0^2,
\qquad
0<c_\kappa\le\frac{2\omega_0^2}{\eta}\Bigl(1+\frac{\rho}{\tau}\Bigr)\le\frac12,
\]
by the definition of \(\omega_0\) in~\eqref{eq:radial-fixed-constants}.

Next, \(\max\{X,|\Lambda|\}=A_0\le\frac14\). Since \(|u+\alpha|\ge\alpha\) and
\(|u+\beta|^2=\beta^2+(1-\beta)\nu\ge\max\{\beta^2,(1-\beta)\nu\}\), we have
\[
|z^0|=|Z-X|=\frac{(1-\alpha)|u|}{|u+\alpha|}X\le\omega_0 X\le\frac18<1,
\qquad
|M|=\frac{\beta}{|u+\beta|}|\Lambda|\le|\Lambda|\le A_0,
\]
and
\[
|\lambda^0|^2=|\Lambda-M|^2
=\frac{(1-\beta)^2\nu}{|u+\beta|^2}|\Lambda|^2
\le(1-\beta)|\Lambda|^2\le A_0^2,
\]
so \(|\lambda^0|\le A_0\le\frac14<1=L_h\). Together with \(x^0=\mu^0=0\), this
gives \(\bu^0\in\mathcal S_0(\mathbf C_0)\).

\paragraph{Step 2: membership in \(\mathcal I(\mathbf C_0,\Delta)\).}
By the radial-Hessian argument in Appendix~\ref{app:lb-rotation-params-proof}
with \(R_*=X\), the function \(f\) in Example~\ref{ex:radial} satisfies
\(L_f\le c_\kappa\le\frac12\) and
\(\|\nabla f(x)\|\le c_\kappa X+b|M|\le\frac32A_0\le\frac38<1\).
Moreover, \(L_{\calA}=0\), \(G_{\calA}=b<1\), and \(h=\|\cdot\|_2\) is convex
and \(1\)-Lipschitz. Since \(|\lambda_c|=|M|\le\frac14\) and
\(\psi(r)\ge-c_\kappa X^2\), every \(x\in\mathbb C\) satisfies
\[
F(x)=\psi(\|x-x_c\|)-b\langle\lambda_c,x-x_c\rangle+b\|x-x_c\|
\ge-c_\kappa X^2+b(1-|\lambda_c|)\|x-x_c\|\ge-c_\kappa X^2.
\]
At the origin, \(\|0-x_c\|=X\), so \(\psi(X)=-\frac{c_\kappa}{2}X^2\) and
\(F(0)\le-\frac{c_\kappa}{2}X^2+b(1+|\lambda_c|)X\le-\frac{c_\kappa}{2}X^2+\frac54A_0\).
Hence
\[
F(0)-\inf F\le\frac{c_\kappa}{2}X^2+\frac54A_0
\le\frac14A_0^2+\frac54A_0<2A_0\le\Delta,
\]
proving \(I\in\mathcal I(\mathbf C_0,\Delta)\).

\paragraph{Step 3: exact rotating trajectory.}
Let \(\mathbb T\) be the linear recursion matrix~\eqref{eq:rot-matrix} on
\(\|x-x_c\|\le R_*\) and \(\mathcal R_\rho\). We verify \(\mathbb T\bv=\zeta\bv\)
directly. Since \(u^2+\nu u+\nu=0\), we have \(u\overline u=\nu\) and
\(\zeta/u=(1+u)/u=-u/\nu\); a direct computation then gives
\[
\frac{u}{u+\alpha}=\frac{u(\overline u+\alpha)}{E_\alpha}=\frac{\nu+\alpha u}{E_\alpha},
\qquad
\frac{\zeta(u+\beta)}{u(u+V)}=-\frac{u(u+\beta)(\overline u+V)}{\nu H}
=\frac{V-\beta}{H}-\frac{\Upsilon}{\nu H}u.
\]
Using \(X-Z=\frac{(1-\alpha)u}{u+\alpha}X\) and the definitions of \(W,b^2,\xi,a\)
in~\eqref{eq:radial-parameters},
\[
u+\chi\frac{(1-\alpha)u}{u+\alpha}-\eta a+\eta\tau b^2\frac{\zeta(u+\beta)}{u(u+V)}
=\Bigl(W-\eta\tau b^2\frac{\Upsilon}{\nu H}\Bigr)u
 +\Bigl(\frac{\nu(W-1)}{\alpha}-\eta a+\xi\Bigr)=0.
\]
Multiplying by \(X\) gives the first row \((\mathbb T\bv)_x=\zeta X\). Next,
\(\Lambda-M=\frac{(1-\beta)u}{u+\beta}\Lambda\) and the definition of \(L\)
give \(u\Lambda+d(\Lambda-M)=\frac{u(u+V)}{u+\beta}\Lambda=\tau b\zeta X\),
verifying the dual row \((\mathbb T\bv)_\lambda=\zeta\Lambda\). The two center
rows follow directly from \((u+\alpha)Z=\alpha\zeta X\) and
\((u+\beta)M=\beta\zeta\Lambda\). Thus \(\mathbb T\bv=\zeta\bv\).

Define \(\widetilde\bu^{\,k}:=\bu_c+\zeta^k\bv\) for \(k\ge0\). Then
\(\widetilde\bu^{\,0}=\bu^0\) and
\(\widetilde\bu^{\,k+1}-\bu_c=\mathbb T(\widetilde\bu^{\,k}-\bu_c)\).
For every \(k\ge0\) and \(j\in\{k,k+1\}\), we have
\(|\widetilde x^{\,j}-x_c|=|\zeta^j X|=X=R_*\) and
\[
|\widetilde\lambda^{\,k}+\rho\calA(\widetilde x^{\,j})|
=|-M+\zeta^k\Lambda+\rho b\zeta^j X|
\le|M|+|\Lambda|+\rho bX
\le(2+\rho)A_0\le1,
\]
so \((\widetilde x^{\,k},\widetilde\lambda^{\,k}),(\widetilde x^{\,k+1},\widetilde\lambda^{\,k})\in\mathcal R_\rho\)
and every primal iterate lies in the quadratic region of \(\psi\). By
Proposition~\ref{prop:embedding} and induction, the actual iterates of
\(\mathrm{LRALM}_{\theta(\epsilon),\alpha,\beta}\) from \(\bu^0\) satisfy
\(\bu^k=\bu_c+\zeta^k\bv\) for all \(k\ge0\).

\paragraph{Step 4: uniform residual lower bound.}
By Proposition~\ref{prop:embedding},
\(R_k\ge|\nabla_\lambda\Psi_\rho(x^{k+1},\lambda^k)|=b|x^{k+1}-x_c|=bX\)
for all \(k\ge0\). Since \(\sqrt H=|u+V|\ge\frac12\sqrt\nu\) and
\(V-\beta\le d<\frac{\omega_0^2}{2}\alpha=\frac{\omega_0}{2}\sqrt\nu\le\sqrt H\),
we have \(|u+\beta|\le|u+V|+(V-\beta)\le2\sqrt H\), and therefore
\[
\frac{b}{|L|}=\frac{\sqrt\nu\sqrt H}{\tau|u+\beta|}\ge\frac{\omega_0\alpha}{2\tau}.
\]
Combining this with \(b\ge\frac{\omega_0\alpha}{\sqrt{\eta\tau}}\) and
\(X=A_0/\max\{1,|L|\}\) gives, for every \(k\ge0\),
\[
R_k\ge bX=A_0\min\left\{b,\;\frac{b}{|L|}\right\}
\ge A_0\omega_0\min\left\{\frac{1}{\sqrt{\eta\tau}},\;\frac{1}{2\tau}\right\}\alpha
=c_\mathrm{r}\alpha>c_\mathrm{r}B\epsilon\ge\epsilon.
\]
Hence \(T_\epsilon(\theta(\epsilon),\alpha,\beta;I,\bu^0)=+\infty\).
\end{proof}

\subsection{Proofs of Theorems~\ref{thm:lower-bound} and~\ref{thm:lower-bound-fixed-rho}}
\label{app:lb-combine-all}

\subsubsection{Proof of Theorem~\ref{thm:lower-bound}}
\label{subsec:lb-combine}

Let \(c_1\) and \(c_2\) be given by~\eqref{eq:rot-thresholds}. The constant
\(c_1\) depends only on \((\gamma,\eta,\tau)\), and \(c_2\) depends only on
\((\gamma,\tau,\Delta)\), as required in Theorem~\ref{thm:lower-bound}.

\begin{proof}[Proof of Theorem~\ref{thm:lower-bound}]
Consider any \(\epsilon\in(0,c_2)\), \(\rho>0\), \(s\ge\tau\) with
\(\epsilon(\rho+s)\le c_1\), and \((\alpha,\beta)\in[0,1]^2\). The
definition of \(c_2\) gives \(\epsilon\le\frac1{64}\le\frac12\),
\(\epsilon\le\sqrt\Delta/2\), and
\(\epsilon^2\le\Delta/(64\gamma)\le\Delta/(16\gamma)\). Note that
\(\zero\in\mathcal S_0(\mathbf C_0)\subseteq\mathcal S_0(\mathbf C)\). We treat
the three parts of the partition of \([0,1]^2\) in Section~\ref{subsec:hard-instances}
separately.

\emph{Case 1: \(\alpha\le\alpha_\epsilon\).} Let \(I\in\mathcal I(\mathbf C_0,\Delta)\)
be the instance of Example~\ref{ex:plateau}. If \(\alpha=0\),
Lemma~\ref{lem:main-plateau} gives
\(T_\epsilon(\theta,\alpha,\beta;I,\zero)=+\infty\). If
\(0<\alpha\le\alpha_\epsilon\), Lemma~\ref{lem:main-plateau} and the
definition of \(\alpha_\epsilon\) give
\[
T_\epsilon(\theta,\alpha,\beta;I,\zero)
>\frac{\Delta}{16\gamma\alpha\epsilon^2}
\ge\frac{\Delta}{16\gamma\alpha_\epsilon\epsilon^2}
=\frac{\Delta\,\epsilon^{-4}}{16384\,\gamma\max\{1,4\tau,4\gamma\tau\}}.
\]
Both bounds hold for every \(\beta\in[0,1]\).

\emph{Case 2: \(\alpha>\alpha_\epsilon\) and \(\beta>0\).} By
Lemma~\ref{lem:main-rotation} (proved via Lemmas~\ref{lem:rot-estimates}
and~\ref{lem:rot-trajectory}), the instance \(I\in\mathcal I(\mathbf C_0,\Delta)\)
of Example~\ref{ex:rotation} satisfies
\(T_\epsilon(\theta,\alpha,\beta;I,\zero)=+\infty\).

\emph{Case 3: \(\alpha>\alpha_\epsilon\) and \(\beta=0\).} By
Lemma~\ref{lem:main-huber}, the instance \(I\in\mathcal I(\mathbf C_0,\Delta)\)
of Example~\ref{ex:huber} satisfies
\(T_\epsilon(\theta,\alpha,0;I,\zero)=+\infty\).

The three cases cover \([0,1]^2\). Since
\(\mathcal I(\mathbf C_0,\Delta)\subseteq\mathcal I(\mathbf C,\Delta)\) and
\(\zero\in\mathcal S_0(\mathbf C)\) for every \(\mathbf C\in\RR^5_{\ge0}\) with
\(L_f,G_f,G_{\calA},L_h\ge1\), taking the supremum over
\((I,\bu^0)\in\mathcal I(\mathbf C,\Delta)\times\mathcal S_0(\mathbf C)\) and
then the infimum over \((\alpha,\beta)\in[0,1]^2\) in
Definition~\ref{def:complexity-measure} gives the claimed lower bound.
\end{proof}

\subsubsection{Proof of Theorem~\ref{thm:lower-bound-fixed-rho}}
\label{subsec:lb-fixed-rho-proof}

\begin{proof}[Proof of Theorem~\ref{thm:lower-bound-fixed-rho}]
With \(B\) and \(\epsilon_\mathrm{r}\) given
in~\eqref{eq:radial-fixed-constants}, set
\(\epsilon_0:=\min\{\epsilon_\mathrm{r},\,\sqrt{\Delta}/2,\,\sqrt{\Delta/(16\gamma)}\}\).
Consider any \(\epsilon\in(0,\epsilon_0)\) and \((\alpha,\beta)\in[0,1]^2\), and
write \(\theta=\theta(\epsilon)\).

\emph{Case 1: \(\alpha\le B\epsilon\).} Since \(\epsilon_\mathrm{r}\le\frac12\),
\(\epsilon\) satisfies the conditions of Example~\ref{ex:plateau} and
Lemma~\ref{lem:main-plateau}. Let \(I\in\mathcal I(\mathbf C_0,\Delta)\) be the
instance of Example~\ref{ex:plateau}. If \(\alpha=0\),
Lemma~\ref{lem:main-plateau} gives
\(T_\epsilon(\theta,\alpha,\beta;I,\zero)=+\infty\). If
\(0<\alpha\le B\epsilon\), it gives
\[
    T_\epsilon(\theta,\alpha,\beta;I,\zero)
    >\frac{\Delta}{16\gamma\alpha\epsilon^2}
    \ge\frac{\Delta\,\epsilon^{-3}}{16\gamma B}.
\]

\emph{Case 2: \(\alpha>B\epsilon\).} Lemma~\ref{lem:main-radial} gives the
instance \(I\in\mathcal I(\mathbf C_0,\Delta)\) and initialization
\(\bu^0\in\mathcal S_0(\mathbf C_0)\) of Example~\ref{ex:radial} with
\(T_\epsilon(\theta,\alpha,\beta;I,\bu^0)=+\infty\).

The two cases cover \([0,1]^2\). Since
\(\mathcal I(\mathbf C_0,\Delta)\subseteq\mathcal I(\mathbf C,\Delta)\) and
\(\mathcal S_0(\mathbf C_0)\subseteq\mathcal S_0(\mathbf C)\) for every
\(\mathbf C\in\RR^5_{\ge0}\) with \(L_f,G_f,G_{\calA},L_h\ge1\), taking the
supremum over \((I,\bu^0)\in\mathcal I(\mathbf C,\Delta)\times\mathcal S_0(\mathbf C)\)
and then the infimum over \((\alpha,\beta)\in[0,1]^2\) in
Definition~\ref{def:complexity-measure} gives the bound.
\end{proof}

